\subsection{\label{subsecPiKP}Charged pions, kaons, and protons}

For the measurement of charged pions, kaons, and protons, several sub-analyses are combined for the comprehensive exploitation of the available PID techniques in ALICE. The spectra cover a range from 0.1/0.2/0.3 GeV/$c$ to 20 GeV/$c$ for \pipm/ \kapm/ \p($\rm\overline{p}$), respectively, with the exact ranges reported in Tab.~\ref{tab:OverviewIndividualAnalysesPiKP}. Similar approaches were followed in earlier analyses in \pp, \pPb\ and \PbPb\ collisions~\cite{Abelev:2013vea,Adam:2015qaa}.  An overview of the individual analyses is presented in Tab.~\ref{tab:OverviewIndividualAnalysesPiKP}. Here, we briefly review the most relevant aspects of previously employed techniques: ITS standalone and TPC-TOF. Additionally, we describe methods which are used for the first time for the measurement of \pt-spectra of charged kaons, pions, and protons: Bayesian PID, TPC-TOF fits, and TPC template fits.

In the ``ITS stand-alone'' technique, the average energy loss in the four outer ITS layers is calculated as a truncated mean. For each particle mass hypothesis, the distance between the measured and the expected value is calculated in multiples of the standard deviation $\sigma$ of the measured energy loss distribution and the particle mass hypothesis with the smallest value is assigned. In contrast to the analysis in the high track density environment of central heavy-ion collisions, the contribution of tracks with wrongly assigned signal clusters is negligible even in the highest pp event multiplicity class. 
In the intermediate \pt-range where a track-by-track identification is feasible, the ``TPC-TOF'' analysis identifies particles by requiring that the measured energy loss signals in the TPC and time-of-flight in the TOF are within 3$\sigma$ of the expected value assuming a specific mass hypothesis. This approach finds its natural limitation towards higher momenta, as the expected energy losses and flight times for different species are insufficiently different to allow for a clear separation. The \pt-ranges in which this procedure is applicable are given in Tab.~\ref{tab:OverviewIndividualAnalysesPiKP}. Two alternative methods, namely Bayesian PID and TPC-TOF fits, were employed in order to unfold the measured d$E$/d$x$ and TOF distributions. The Bayesian method of particle identification for the extraction of the minimum-bias spectra of pions, kaons, and protons is described in detail in~\cite{Adam:2016acv}. The {\it a priori} probabilities used in the Bayesian-approach analysis were extracted from the experimental data for the minimum bias event sample using an iterative procedure. The influence of different sets of {\it a priori} probabilities, determined for the lowest and highest event multiplicity bins, was evaluated and included in the systematic uncertainties. The actual identification of particles is based on the maximum probability method in which the most likely particle type is assigned to the track. 

\begin{table}[tbph!]
        \centering
        \begin{tabular}{cccccc}
\hline\hline
                Analysis & PID & \multicolumn{3}{c}{\pt\ range (\GeVc)} & (pseudo)rapidity\\
                \cline{3-5}
                         & technique                                                                         & \pipm            & \kapm         & p(\pbar)       & range\\ 
\hline
                \multirow{2}{*}{ITS stand-alone}  & n-$\sigma$ cuts  & \multirow{2}{*}{$0.1-0.6$}     & \multirow{2}{*}{$0.2-0.6$}   & \multirow{2}{*}{$0.3-0.6$}   & \multirow{2}{*}{$|y| < 0.5$}  \\ 
                                                  & on ITS & & & & \\
                                           &                                                        &                      &                   &                   &                    \\
                \multirow{2}{*}{Bayesian PID} & Bayesian & \multirow{2}{*}{$0.2-2.5$} & \multirow{2}{*}{$0.3-2.5$}   & \multirow{2}{*}{$0.5-2.5$}   & \multirow{2}{*}{$|y| < 0.5$}  \\
                                           & probability & & & & \\
                                           &                                                        &                      &                   &                   &                    \\
                \multirow{2}{*}{TPC-TOF}   & n-$\sigma$ cuts on & \multirow{2}{*}{$0.25-1.2$} & \multirow{2}{*}{$0.3-1.2$} & \multirow{2}{*}{$0.45-2.0$} & \multirow{2}{*}{$|y| < 0.5$}  \\
                                           & TPC and TOF & & & & \\
                                           &                                                        &                      &                   &                   &                    \\
                \multirow{2}{*}{TPC-TOF fits} & n-$\sigma$ fits    & \multirow{2}{*}{$0.25-2.5$}   & \multirow{2}{*}{$0.3-2.5$}   & \multirow{2}{*}{$0.45-2.7$} & $|y| < 0.5$ (TPC)\\
                                           &  to TPC and TOF     &  &  &  &$|\eta| < 0.2$ (TOF)\\
                                           &                                                        &                      &                   &                   &                    \\
                TPC template & TPC \dEdx        & \multicolumn{3}{c}{\multirow{2}{*}{$> 2.0$}}         & \multirow{2}{*}{$|\eta| < 0.8$}\\
                Fits   &   Template fits   & \multicolumn{3}{c}{} & \\
\hline\hline
        \end{tabular}
        ~\newline
                \caption{Overview of \pt\ ranges used for the combination of the various techniques used for identifying pions, kaons and protons. Since the true rapidity is not known at reconstruction level, fit based analyses (``TPC template fits" and ``TPC-TOF fits"), which determine the yield of pions, kaons and protons simultaneously, require an additional $\eta$-cut. }
        \label{tab:OverviewIndividualAnalysesPiKP}
\end{table}

In the ``TPC-TOF Fits'' method, the energy loss distribution in the TPC is simultaneously fitted by three gaussian distributions corresponding to charged pions, kaons and protons in each \pt\ and multiplicity bin. Similarly, the velocity distribution of the TOF is fitted for all three species simultaneously. In order to guarantee a sufficient separation of the particle species by minimizing the difference between total and transverse momentum \pt, the TOF fits were performed in a narrow $\eta$-window ($|\eta|<0.2$) and afterwards transformed to the common rapidity window of $|y|<0.5$ assuming a flat distribution in $y$.
Above a \pt\ of $\sim2$~GeV/$c$, particle identification can still be achieved statistically, rather than on a track-by-track basis, by fitting the specific energy loss in the relativistic rise region with a multi-component fit function, as done in the ``TPC Template fits'' approach. In this method, the measured d$E$/d$x$ distribution, in which the distributions of several particle species are overlapping, is fitted with a sum of templates (one for each particle type). The templates are extracted in a data-driven procedure from a pure sample of tagged particles of a given type. This pure sample is obtained from weak decay daughter tracks (\p\ and \pbar\ from $\lmb$ and $\almb$ as well as \pipm\ from \kzero) and tracks identified with the TOF (\pipm, \kapm, \p, and $\rm\overline{p}$). After a further strict selection of primary-particle-like topologies, the expected d$E$/d$x$ response is determined in fine bins of momentum and pseudorapidity. The template for each particle species in a given transverse momentum bin in the rapidity window $|y|<0.5$ is then obtained by sampling the measured momenta and pseudorapidity values of the tracks in this bin. The individual particle yields are the only free parameters in the fit of the templates to the measured d$E$/d$x$ distribution.


For all particle species and sub-analyses, contamination from secondary particles at low transverse momenta was subtracted in a data-driven approach on the basis of the measured distance of closest approach of the track to the primary vertex in the transverse plane (${\rm DCA_{xy}}$), as done in previous work~\cite{Adam:2015qaa}. The ${\rm DCA_{xy}}$ distribution of the selected tracks was fitted with three Monte Carlo templates corresponding to the expected shapes of primary particles, of secondaries from material (including electrons from photon conversions) and of secondaries from weak decays. The procedure was repeated for each \pt\ and event multiplicity bin and thus takes into account possible differences in the feed-down correction due to a change of the abundances and spectral shapes of weakly decaying strange particles. 

The efficiencies obtained for anti-protons and kaons have been additionally corrected based on a comparison of the absorption cross-section used in GEANT3 and the more realistic description of hadronic cross-sections in FLUKA, as in~\cite{Adam:2015qaa,Abbas:2013rua}. 

The determination of systematic uncertainties follows the procedures established in previous analyses~\cite{Abelev:2013vea,Adam:2015qaa}. All the considered contributions are summarized in Tab.~\ref{tab:syspikp}. Corrections for secondary particles lead to uncertainties of up to 4\% for protons and 1\% for pions while they are negligible for kaons. The uncertainty in the material budget is of 5\% at very low momenta and is related to the energy loss of the particles in the detector material. In addition, inelastic and elastic hadronic scattering processes inside the detector material are described by the transport codes only with limited precision and lead to uncertainties of up to 6\% for \pbar\ (for which the respective cross-section is largest) at low transverse momenta. The track quality selection criteria and the matching of TPC tracks with ITS hits give rise to a systematic uncertainty of the global tracking efficiency that amounts to 4\%, independent of \pt\ and particle species. 
The Lorentz force causes shifts of the cluster position in the ITS, pushing the charge in opposite directions depending on the polarity of the magnetic field of the experiment ($E \times B$ effect). In the ITS stand-alone analysis, the uncertainty related to this effect is estimated by analysing data samples collected with opposite magnetic field polarities, for which a difference of 3\% is observed.
For those sub-analyses (Bayesian PID, TPC-TOF, TPC-TOF fits) that require in addition that the track under study is matched to a hit in the TOF, an additional uncertainty of 3\%/6\%/4\% is taken into account for pions, kaons, and protons, respectively. Following the approach presented in~\cite{Abelev:2013vea}, this matching efficiency uncertainty was estimated by repeating the analysis separately for those regions in azimuth in which modules of the Transition Radiation Detector were already present in 2010 and for those in which they were not yet installed.


\begin{table*}[ph!]
  \begin{tabularx}{\textwidth}{p{5.3cm}*{3}{Y}*{3}{Y}*{3}{Y}}
\hline\hline
    &  \multicolumn{9}{c}{Uncertainty (\%)} \\
    {\bf Common source}  & \multicolumn{3}{c}{\pipm} & \multicolumn{3}{c}{\kapm} & \multicolumn{3}{c}{\p\ (\pbar)} \\ 
\hline
    \pt\ (\GeVc)   &   0.1 & 3.0 & 20.0   &   0.2 & 2.5 & 20.0 &   0.3 & 4.0 & 20.0   \\
    \cmidrule(l{5pt}){2-4} \cmidrule(l{5pt}){5-7} \cmidrule(l{5pt}){8-10} 
%%% Common Uncertainties 
% Not obvious: would we prefer quoting only 2 and not 3 pt ranges here? 
% --- it might depend on the pT dendence of these systematics 
    Correction for secondaries & 1  & 1 & 1 & \multicolumn{3}{c}{negl.} & 4 & 1 & 1 \\
    Material budget & 5 & \multicolumn{2}{c}{negl.} & 2 & \multicolumn{2}{c}{negl.} & 4 & \multicolumn{2}{c}{negl.} \\
    Hadronic interactions & 2 & 1 & 1 & 3 & 1 & 1 & 4(6) & 1(1) & 1(1) \\
    Global tracking efficiency \newline (incl. track cut variation) & \multicolumn{3}{c}{\multirow{2}{*}{4}} & \multicolumn{3}{c}{\multirow{2}{*}{4}} & \multicolumn{3}{c}{\multirow{2}{*}{4}} \\
    TOF matching efficiency \newline (Bayes.,TPC-TOF,TPC-TOF fits) & \multicolumn{3}{c}{\multirow{2}{*}{3}} & \multicolumn{3}{c}{\multirow{2}{*}{6}} & \multicolumn{3}{c}{\multirow{2}{*}{4}} \\
\hline
    {\bf Specific source} & \multicolumn{3}{c}{\pipm} & \multicolumn{3}{c}{\kapm} & \multicolumn{3}{c}{\p\ (\pbar)} \\ 
\hline
%%% ITS Standalone Analysis Uncertainties 
    \pt\ (\GeVc)   &   0.1 & & 0.6   &   0.2 & & 0.6 &   0.3 & & 0.6   \\
    \cmidrule(l{5pt}){2-4} \cmidrule(l{5pt}){5-7} \cmidrule(l{5pt}){8-10} 
    ITSsa tracking efficiency & 3 &  & 3 & 3 &  & 3 & 3 &  & 3  \\ 
    $E \times B$ effect & \multicolumn{9}{c}{3} \\
    ITS PID & 5 & & 1 & 5 & & 9 & 8 & & 6 \\ 
\hline\hline

    \pt\ (\GeVc)   &   0.2 &  & 2.5   &   0.3 & & 2.5  &   0.5 & & 2.5    \\
    \cmidrule(l{5pt}){2-4} \cmidrule(l{5pt}){5-7} \cmidrule(l{5pt}){8-10} 
    Bayesian PID  & \multicolumn{3}{c}{1} & 1 & & 3 & 1 & & 2 \\   %\multicolumn{3}{c}{1.5\%} & \multicolumn{3}{c}{3.5\%} & \multicolumn{3}{c}{2.5\%} \\
\hline\hline
    
    \pt\ (\GeVc)   &   0.25 &  & 1.2   &   0.3 & & 1.2  &   0.45 & & 2.0    \\
    \cmidrule(l{5pt}){2-4} \cmidrule(l{5pt}){5-7} \cmidrule(l{5pt}){8-10} 
    TPC-TOF PID & \multicolumn{3}{c}{1} & 1 & & 5 & \multicolumn{3}{c}{1} \\
\hline\hline
    
    \pt\ (\GeVc)   &   0.25 &  & 2.5   &   0.3 & & 2.5  &   0.45 & & 2.7    \\
    \cmidrule(l{5pt}){2-4} \cmidrule(l{5pt}){5-7} \cmidrule(l{5pt}){8-10} 
    TPC-TOF fits PID  & 1 & & 5 & 1 & & 10 & 1 & & 8 \\
\hline\hline
    
    \pt\ (\GeVc)   &  2.0 &  & 20.0   &   2.0 & & 20.0  &   2.0 & & 20.0    \\
    \cmidrule(l{5pt}){2-4} \cmidrule(l{5pt}){5-7} \cmidrule(l{5pt}){8-10} 
    TPC template fits PID & 4 & & 6 & 10 & & 12  & 8 & & 13 \\
\hline\hline
   

   {\bf Total}  & \multicolumn{3}{c}{\pipm} & \multicolumn{3}{c}{\kapm} & \multicolumn{3}{c}{\p\ (\pbar)} \\ 
\hline
    \pt\ (\GeVc)   &   0.1 & 3.0 & 20.0   &   0.2 & 2.5 & 20.0 &   0.3 & 4.0 & 20.0   \\
    \cmidrule(l{5pt}){2-4} \cmidrule(l{5pt}){5-7} \cmidrule(l{5pt}){8-10} 
    Total & 8.4 & 5.0 & 7.2 & 7.5 & 6.6 & 12.6 & 12.3 & 15.1 & 13.3 \\
    Total ($N_{ch}$-independent) & 8.1 & 4.4 & 6.9 & 6.7 & 6.1 & 12.2 & 10.5 & 13.5 & 11.5\\
    \bottomrule
  \end{tabularx}
  ~\newline
    \caption{Main sources and values of the relative systematic uncertainties of the \pt-differential yields of \pipm, \kapm\ and \p (\pbar). The values are reported for low, intermediate and high \pt. The contributions that act differently in the various event classes are removed from the Total (quadratic sum of all contributions), defining the $N_{ch}$-independent ones, which are correlated across different multiplicity intervals. The contribution from the global tracking efficiency is common to all analyses except for the ITS stand-alone (ITSsa). 
}
  \label{tab:syspikp}  
\end{table*}


All sub-analyses were found to be in agreement in the overlapping \pt\ ranges within the uncorrelated part of their respective systematic uncertainties. The final combined spectrum for each particle species was then obtained by calculating the average over all sub-analyses using the uncorrelated part of their systematic errors as weights~\cite{Patrignani:2016xqp}. 
The uncertainties originating from common sources were then added in quadrature to each other and to the uncertainty attributed to the specific particle identification methods. The systematic uncertainty contribution that is uncorrelated across multiplicities was estimated to be $\sim$4-8\%, $\sim$6-12\% and $\sim$10-14\% for \pipm, \kapm and \p, respectively, for all V0M multiplicity bins. 


